\documentclass[11pt,letterpaper]{article}

% Enable the inclusion of EPS-format graphics.
\usepackage{epsfig}
% Used to handle line spacing. Must be loaded before hyperref.
\usepackage{setspace}
% Clickable links.
\usepackage[bookmarks,colorlinks=true,linkcolor=blue,urlcolor=blue]{hyperref}
% Code listings via \lstlisting.
\usepackage{listings}
% Multirow table cell spanning.
\usepackage{multirow}
% Zero paragraph indent and non-zero paragraph separation.
\usepackage{parskip}
% Tables with flexible-width columns.
\usepackage{tabulary}
% Diagrams and stuff.
\usepackage{tikz}

% Set document fonts: Times as the default Roman font, Helvetica for
% sans-serif, Courier for teletype.
\renewcommand{\rmdefault}{ptm}
\renewcommand{\sfdefault}{phv}
\renewcommand{\ttdefault}{pcr}

% I want captions to have slightly smaller text and a bold label.
\usepackage[font=small,labelfont=bf]{caption}

% Customize formatting for code listings.
\lstset{
    basicstyle={\small\ttfamily}
}

% LaTeX' default line spacing is a little tight for my liking.
\setstretch{1.15}

\begin{document}

\title{\includegraphics[width=0.75\textwidth]{logo.eps}\\
Problem Definition Format}
\author{Samuel Coleman}
\date{\today}

% Turn off index link generation for the unnumbered title page so we don't get
% a warning about a missing page.
\hypersetup{pageanchor=false}
\begin{titlepage}

\maketitle
% Ditch the page styling (including page number) for the title page.
\thispagestyle{empty}

\end{titlepage}
% Having generated the page, we can turn index link generation back on.
\hypersetup{pageanchor=true}

\tableofcontents

\newpage

\section{Introduction}
\label{introduction}

In lieu of the absence of a standard format for distributing problems sets for
online judges, Garrit proposes that components implementing the rest of its
interfaces seek to maintain compatibility with the problem format used by its
reference components.

Problem sets are considered to have three components: the description, which
educates the user about the problem being solved and introduces the constraints
and expectations of the solution; sample data or ``samples'', which help the
user better understand the definition and implement their solution; and judging
data or ``cases'', against which user submissions are judged. The latter two
components consist primarily of input data to be fed into submissions, and
output data against which the submitted program output is to be compared.

\section{Organization}
\label{organization}

Each problem set should be stored in its own directory, named as the problem.
At minimum, the directory must contain the problem definition as a JSON
document in a file named \texttt{problem.json}, described below. It is strongly
recommended that the directory also contain a ``readme'' file, and that the
problem definition leverage the contents of this file as the problem
description.

Sample problem definitions are available as part of
Garrit\footnote{\url{https://github.com/Garrit/problems}}. As an example, here
is the directory structure for the ``echo'' example:

\begin{tabular}{ l l }
    \texttt{echo/} & Root directory \\
    \texttt{\hphantom{echo/}problem.json} & Problem definition \\
    \texttt{\hphantom{echo/}README.md} & Problem description \\
    \texttt{\hphantom{echo/}samples/} & Sample input/output data \\
    \texttt{\hphantom{echo/samples/}sample.in} & \hphantom{Sample} input \\
    \texttt{\hphantom{echo/samples/}sample.out} & \hphantom{Sample} output \\
    \texttt{\hphantom{echo/}cases/} & Secret judging data \\
    \texttt{\hphantom{echo/cases/}case1.in} & \hphantom{Secret} input \#1 \\
    \texttt{\hphantom{echo/cases/}case1.out} & \hphantom{Secret} output \#1 \\
    \texttt{\hphantom{echo/cases/}case2.in} & \hphantom{Secret} input \#2 \\
    \texttt{\hphantom{echo/cases/}case2.out} & \hphantom{Secret} output \#2 \\
\end{tabular}

This structure is, while consistent throughout the sample problems, entirely
arbitrary and controlled by the problem definition. While it might be nice for
others using your problem sets if they follow a logical ordering, no structure
beyond the placement of \texttt{problem.json} is enforced.

\section{Problem Definition Format}
\label{format}

The file \texttt{problem.json} must contain a JSON object with the following
properties:

\nopagebreak
\begin{tabulary}{\textwidth}{ | l | l | L | }
    \hline
    \textbf{Key} & \textbf{Type} & \textbf{Value} \\
    \hline
    \texttt{name} & string & The name of the problem set. Should be consistent
        with the name of the directory in which the problem definition is
        stored. \\
    \hline
    \texttt{description} & string & The problem description. We suggest the use
        of
        Markdown\footnote{\url{http://daringfireball.net/projects/markdown/}}
        or a derivative syntax for formatting.
        \newline
        \emph{May not be present if the \texttt{descriptionFile} property is
        present.} \\
    \hline
    \texttt{timeLimit} & integer & A default limit in seconds for the execution
        time of each case. If not specified, the execution environment may
        impose its own limit. \emph{Optional.} \\
    \hline
    \texttt{memoryLimit} & integer & A default limit in bytes for the memory
        consumption of each case. If not specified, the execution environment
        may impose its own limit. \emph{Optional.} \\
    \hline
    \texttt{descriptionFile} & string & The name of a file within the problem
        set directory containing the problem description. As with descriptions
        embedded within the problem definition, Markdown formatting syntax is
        suggested.
        \newline
        \emph{May not be present if the \texttt{description} property is
        present.} \\
    \hline
    \texttt{samples} & array & An array of objects containing sample data for
        the user. \\
    \hline
    \texttt{cases} & array & An array of objects containing the secret judging
        data. \\
    \hline
\end{tabulary}

Each object describing a sample should have the following properties:

\nopagebreak
\begin{tabulary}{\textwidth}{ | l | l | L | }
    \hline
    \textbf{Key} & \textbf{Type} & \textbf{Value} \\
    \hline
    \texttt{name} & string & A name identifying the data set. It should be
        unique within the problem set. \\
    \hline
    \texttt{input} & string & The problem input data. As with the
        \texttt{description} and \texttt{descriptionFile} properties of the
        top-level problem set object, the input data may be either provided
        within the problem definition, or held in an external file.
        \newline
        \emph{May not be present if the \texttt{inputFile} property is
        present.} \\
    \hline
    \texttt{inputFile} & string & The name of a file within the problem
        set directory containing the sample input.
        \newline
        \emph{May not be present if the \texttt{input} property is present.} \\
    \hline
    \texttt{output} & string & Expected output of a solution for the
        corresponding input.
        \newline
        \emph{May not be present if the \texttt{outputFile} property is
        present.} \\
    \hline
    \texttt{outputFile} & string & The name of a file within the problem
        set directory containing the sample output.
        \newline
        \emph{May not be present if the \texttt{output} property is present.}
        \\
    \hline
\end{tabulary}

Objects describing cases are identical to those describing samples, but may
optionally have extra parameters:

\nopagebreak
\begin{tabulary}{\textwidth}{ | l | l | L | }
    \hline
    \textbf{Key} & \textbf{Type} & \textbf{Value} \\
    \hline
    \texttt{timeLimit} & integer & A limit in seconds for the execution time of
        the case. Overrides any default limit specified in the problem
        definition. \emph{Optional.} \\
    \hline
    \texttt{memoryLimit} & integer & A limit in bytes for the memory
        consumption of the case. Overrides any default limit specified in the
        problem definition. \emph{Optional.} \\
    \hline
\end{tabulary}

\end{document}